\section{A path bridge under the exact physical singleton}
\label{sec:v36-exact-path-bridge}
We retain the original stationary observation
$A_{t,\lambda}(k,m)=\{W_\lambda(t)=k,C_\lambda(t)=m\}$.
Its conditional path law has a Brownian bridge limit. The proof needs
both a local finite-dimensional calculation and the preceding local
fourth moments. A total-variation approximation of the posterior by
itself would not identify this bridge.

For a positive definite $3\times3$ matrix $V$ and $\xi\in\R^3$,
write $\mathsf{Br}_{V,\xi}$ for the law of
\begin{equation}\label{eq:v36-bridge-definition}
 s\xi+V^{1/2}(B(s)-sB(1)),\qquad 0\le s\le1,
\end{equation}
where $B$ is standard three-dimensional Brownian motion. In particular
its covariance at $(s,u)$ is $(\min\{s,u\}-su)V$. We use the
bounded-Lipschitz metric on continuous path space with the supremum
norm. This is weak convergence of path laws, not convergence in total
variation to a Gaussian measure.

\subsection{Fourier probes of finitely many collision blocks}
Fix $0=s_0<s_1<\cdots<s_q=1$, put
$m_j=\lfloor ms_j\rfloor-\lfloor ms_{j-1}\rfloor$, and choose fixed
real vectors $v_1,\ldots,v_q$. Let
\[
 \mathcal P_m(x)=\exp\left\{i\sum_{j=1}^q
       v_j\cdot\frac{V_{\lfloor ms_{j-1}\rfloor,m_j}(x)}{\sqrt m}\right\}.
\]
This multiplier has modulus one on the original collision trajectory.

\begin{lemma}[Local block probes and endpoint measures]
\label{lem:v36-local-block-probes}
Along $\lambda_m\to\lambda_0$ and
\[
 \zeta_m=\frac{(k_m,t_m-m\bar\tau_{\lambda_m})}{\sqrt m}
                  \longrightarrow\zeta,
\]
the complex endpoint measures
\begin{equation}\label{eq:v36-probed-endpoint-measure}
 m^{3/2}(x,S_m-t_m,T_{\lambda_m}^mx)_*
              (\mathcal P_m\one_{\{K_m=k_m\}}\nu)
\end{equation}
converge vaguely, with uniformly bounded variation on bounded time slabs, to
\[
 g_{\Sigma_{\lambda_0}}(\zeta)\,
 \mathcal C_{\Sigma_{\lambda_0},\zeta}(\mathbf v)\,
                      \nu\otimes ds\otimes\nu,
\]
where, with $\alpha_j=s_j-s_{j-1}$ and $\bar v=\sum_j\alpha_jv_j$,
\begin{equation}\label{eq:v36-bridge-characteristic}
 \mathcal C_{\Sigma,\zeta}(\mathbf v)
 =\exp\left\{i\bar v\cdot\zeta
       -\tfrac12\left(\sum_j\alpha_jv_j^{\mathsf T}\Sigma v_j
                                  -\bar v^{\mathsf T}\Sigma\bar v\right)\right\}.
\end{equation}
This is the characteristic function of consecutive increments of
$\mathsf{Br}_{\Sigma,\zeta}$. The convergence allows the varying
endpoint-overlap tests of \textup{(H7)}.
\end{lemma}
\begin{proof}
First use smooth endpoint functions $a,d$ and a time test $q_0$ whose
Fourier transform is integrable and has fixed compact support.
The exact Fourier integral contains the chronological word
\begin{equation}\label{eq:v36-probe-product}
 \ell\left(M_d Q_{\omega+v_q/\sqrt m}^{m_q}\cdots
               Q_{\omega+v_1/\sqrt m}^{m_1}(a\nu)\right)
\end{equation}
and the factor
$e^{-i[u\cdot k_m+b(t_m-m\bar\tau)]}\widehat q_0(-b)$.
Every $m_j$ is a positive fixed fraction of $m$ for all sufficiently
large $m$. On the complement of a fixed neighborhood of zero the
shifted frequencies remain away from zero, so the product is
exponentially small by the fixed-band spectrum. Near zero put
$w=\sqrt m\,\omega$ and use the analytic splitting in every block.
The product of projections tends to the unperturbed rank-one
projection; its endpoint amplitude is $\nu(a)\nu(d)$.
After taking out the factor $\nu(a)\nu(d)\widehat q_0(0)$,
the remaining Gaussian inverse is
\[
 \frac1{(2\pi)^3}\int e^{-iw\cdot\zeta}
      \exp\left\{-\tfrac12\sum_j\alpha_j
                        (w+v_j)^{\mathsf T}\Sigma(w+v_j)\right\}\,dw.
\]
Its value is $g_\Sigma(\zeta)\mathcal C_{\Sigma,\zeta}(\mathbf v)$.
This follows by the elementary Fourier transform of a Gaussian,
or by conditioning independent Gaussian increments of covariances
$\alpha_j\Sigma$ on their sum. The identity also fixes the sign of
the mean in \eqref{eq:v36-bridge-characteristic}. A bound of the form
$C_{\mathbf v}e^{-c|w|^2}$ dominates the rescaled product, uniformly
on the local parameter cover. The cubic remainders vanish for each
fixed $w$, so dominated convergence is legitimate.

To pass from $q_0$ to an interval $J$, take the fixed-band upper and
lower envelopes $q_B^+\ge\one_J\ge q_B^-$. For nonnegative smooth
endpoints, the absolute error in replacing the probed interval integral
by its upper-envelope integral is bounded by the unprobed integral
with the positive test $q_B^+-q_B^-$. The latter has limiting mass
$4\pi/B$ times the endpoint Gaussian amplitude. Take $m\to\infty$
with $B$ fixed and then $B\to\infty$. This is a domination argument
for the complex probe, not an order comparison of complex numbers.
Continuous endpoint functions follow by uniform approximation and
decomposition into real and imaginary parts.

The total variation of \eqref{eq:v36-probed-endpoint-measure} is
bounded by the unprobed positive endpoint measure. Its masses on
bounded time slabs are uniformly bounded. Products of endpoint
functions and interval tests therefore determine convergence for
continuous compactly supported tests. The same domination extends it
to tests continuous off a limiting null set. For varying tests remove
an exceptional open neighborhood with small-measure closure and null
boundary. Portmanteau for the dominating positive measures controls
that closure; uniform convergence controls the complement. This is
exactly the condition in \textup{(H7)}. It proves the last assertion
without a new regularity requirement on a long pulled-back probe.
\end{proof}

\subsection{The original physical-time bridge}
Let $\mathcal X_{t,\lambda}$ be a continuous interpolation of
\begin{equation}\label{eq:v36-physical-path}
 t^{-1/2}\bigl(W_\lambda(ts),C_\lambda(ts)-ts/\bar\tau_\lambda\bigr),
                         \qquad 0\le s\le1.
\end{equation}
It can be obtained by joining its values at collision and cell-crossing
times, and at $0,t$. Since all one-flight label changes and collision
jumps are uniformly bounded, its supremum distance from the step
version is $O(t^{-1/2})$. This also defines the same limiting law in
Skorokhod space for the original step path.

\begin{theorem}[Microscopic conditional path bridge]
\label{thm:v36-microscopic-path-bridge}
For every compact family satisfying \textup{(H1)--(H7)} and each fixed
$M<\infty$,
\begin{align}
 \sup_{\substack{\lambda\in\mathcal K,
                     |\xi_{t,\lambda}(k,m)|\le M}}
 d_{\rm BL}\left(
  \operatorname{Law}_{\mu_\lambda}(\mathcal X_{t,\lambda}
                              \mid A_{t,\lambda}(k,m)),
  \mathsf{Br}_{\mathcal V_\lambda,\xi_{t,\lambda}(k,m)}\right)
                          \longrightarrow0.
 \label{eq:v36-exact-physical-bridge}
\end{align}
The probability in the conditioning is the original singleton
probability, positive for all sufficiently large $t$ on that set.
The theorem remains valid under nonnegative bounded regular initial
and terminal selectors satisfying \textup{(H7)} and the uniform
positive product-of-means condition.
\end{theorem}
\begin{proof}
Consider first any convergent sequence of parameters and central
endpoints, with $\xi_{t,\lambda}(k,m)\to\xi$. It has
$m/t\to1/\bar\tau_{\lambda_0}$ and a bounded collision endpoint
$\zeta_m=(k,t-m\bar\tau_\lambda)/\sqrt m$.
Pass to a subsequence along which $\zeta_m\to\zeta$.

Integrate the exact flight age in the numerator with a block probe.
For each pair of cell offsets $(d,e)$, the collision label is
$k+d-e$, and the overlap is
\[
 F_\lambda^{d,e}(x,s,y)=
           \int\one_{I_{\lambda,d}(x)}(a)
                       \one_{I_{\lambda,e}(y)}(a-s)\,da,
                        \qquad s=S_m-t.
\]
Apply Lemma~\ref{lem:v36-local-block-probes}. The finitely many label
shifts change the rescaled endpoint by $O(m^{-1/2})$. Consequently
every term has the same limiting bridge characteristic function.
The sum of its overlap amplitudes is $\bar\tau_\lambda^2$, and the
stationary source supplies $1/\bar\tau_\lambda$. Dividing by the
already proved positive local denominator leaves precisely
\eqref{eq:v36-bridge-characteristic}. Thus the collision path $Z_m$
has the finite-dimensional conditional limits of
$\mathsf{Br}_{\Sigma_{\lambda_0},\zeta}$.
Corollary~\ref{cor:v36-conditional-tightness} proves tightness in the
continuous path space. This identifies the full conditional collision
path law. No independence of adjacent deterministic collision blocks
was assumed: their asymptotic calculation came from
\eqref{eq:v36-probe-product}.

It remains to justify the physical clock under the same conditioning.
Let $a$ be the original initial flight age and let $a_s$ be the age
at physical time $ts$. On the conditioned event, set
$r_s=C_\lambda(ts)$; then $0\le r_s\le m$ and
\begin{equation}\label{eq:v36-bridge-exact-clock}
 S_{r_s}=ts+a-a_s,\qquad
 \mathcal X_{t,\lambda}(s)
   =\sqrt{m/t}\,D_\lambda^{\rm clk} Z_m(r_s/m)+O(t^{-1/2})
\end{equation}
uniformly in $s$, allowing the stated interpolation. The spatial
error in the second identity is the difference of the current and
initial within-flight cell labels. Both labels belong to the same
fixed finite set used in the overlap formula. Thus neither endpoint
cell correction has been lost.

Conditional tightness of $Z_m$ implies
$\max_{j\le m}|S_j-j\bar\tau_\lambda|=O_P(\sqrt m)$ under these
conditional laws, uniformly over the compact central set. The first
identity, $0\le a,a_s\le\tau_+$, and
$t-m\bar\tau_\lambda=O(\sqrt m)$ give
$\sup_s|r_s/m-s|\to0$ in conditional probability. Tightness with
continuous limits then permits replacement of $r_s/m$ by $s$ in
\eqref{eq:v36-bridge-exact-clock}. Independence of the clock and the
collision path is not needed.

The limiting linear image has covariance
$\bar\tau_{\lambda_0}^{-1}D_{\lambda_0}^{\rm clk}
\Sigma_{\lambda_0}(D_{\lambda_0}^{\rm clk})^{\mathsf T}
=\mathcal V_{\lambda_0}$. Its endpoint is correct because the exact
algebraic identity
\[
 \sqrt{m/t}\,D_\lambda^{\rm clk}
        \frac{(k,t-m\bar\tau_\lambda)}{\sqrt m}
                         =\xi_{t,\lambda}(k,m)
\]
holds. Therefore the limit is
$\mathsf{Br}_{\mathcal V_{\lambda_0},\xi}$.
Compactness, uniform tightness, and continuity of the Gaussian bridge
law in $(\mathcal V,\xi)$ turn this sequential conclusion into the
uniform distance in \eqref{eq:v36-exact-physical-bridge}.

With regular endpoint selectors the age overlap is weighted before
integration. Its limiting amplitude is the product of stationary
selector means times the same normalization as in the selected local
law. This positive amplitude cancels with the selected denominator.
The conditional fourth moments are covered by the selected part of
Corollary~\ref{cor:v36-conditional-tightness}. All the remaining clock
identities concern the same physical trajectory and are unchanged.
\end{proof}

\begin{corollary}[The positive-band posterior has the same bridge]
\label{cor:v36-band-posterior-bridge}
Under the conditions of Theorem~\ref{thm:v36-microscopic-path-bridge},
let the posterior be formed with the positive flight-age likelihood
of Corollary~\ref{cor:v35-family-posterior}, at
$B_m=m^{P+3/2}$ for fixed $P>0$. Its pushforward by the original
physical path map has the same bridge limit. Its total-variation
norm distance from the exact conditional path law is $O(m^{-P})$.
The distance of either law to the Gaussian bridge is only asserted
to tend to zero.
\end{corollary}
\begin{proof}
Pushforward contracts the total-variation norm. Apply the proved
source-posterior comparison and then
Theorem~\ref{thm:v36-microscopic-path-bridge}. The polynomial rate
belongs to comparison of two posteriors on the same source, not to
the weak bridge limit. All spectral bands in the proof of that limit
were fixed before passing to large $m$.
\end{proof}

\paragraph{Relation to the return-density problem.}
The new theorem closes the microscopic bridge for the stationary
physical singleton, including its regular endpoint-selected form.
It is stronger than the inherited window bridges because its
conditioning is the exact event, with probability of order $t^{-3/2}$.
It neither supplies an arbitrary path-selector Gaussian amplitude nor
controls the common pointwise correction or full frequency complement
for the four-coordinate record $J_{n,R}$. Those quantities in Part II
retain their original definitions. No return event was substituted
for a physical observation in this proof.
